% =============================================================================
% Chapter 6 --- Monitoring, Alerting, and Audit Schema
% =============================================================================
\chapter{Monitoring, Alerting, and Audit Schema}
\label{ch:monitoring}

\section{Monitoring objectives}

Post-deployment monitoring serves four objectives:

\begin{enumerate}
  \item \textbf{Detect policy violations and drift} in real-time
  \item \textbf{Maintain complete audit trail} for governance evidence
  \item \textbf{Enable rapid incident response} and containment
  \item \textbf{Support continuous improvement} through metrics analysis
\end{enumerate}

\section{Required audit schema}

Every governed execution must record the following fields per
\filepath{docs/schema/regulations/audit-event.schema.json}:

\begin{table}[htbp]
\centering
\footnotesize
\caption{Required Audit Event Fields}
\label{tab:audit-fields}
\begin{tabularx}{\textwidth}{@{}l l X@{}}
\toprule
\textbf{Field} & \textbf{Type} & \textbf{Description} \\
\midrule
\texttt{event\_id} & string & Unique audit event identifier (e.g., evt-20260422-abc123) \\
\texttt{event\_type} & enum & Discriminator: agent\_inference, agent\_write, review\_requested, review\_approved, review\_rejected, agent\_error \\
\texttt{timestamp} & datetime & ISO 8601 timestamp when event was recorded \\
\texttt{request\_identity} & object & Request ID, correlation ID, user ID, tenant ID \\
\texttt{agent\_identity} & object & Agent ID, type, version, service name, deployment ID \\
\texttt{action} & object & Action type, entity ID, input/output hashes \\
\texttt{approval\_chain} & array & Ordered approval steps if HITL invoked \\
\texttt{outcome} & object & Status (success/failure/pending), duration\_ms, error\_code \\
\bottomrule
\end{tabularx}
\end{table}

\section{Identity attribution requirements}

Per \controlid{AGT-03} and REG-MGF-2.1.2-002, every action must be attributable:

\subsection{Agent identity fields}

\begin{itemize}
  \item \texttt{agent\_id}: Unique identifier for the agent instance
  \item \texttt{agent\_type}: Classification (e.g., predictive-analytics)
  \item \texttt{agent\_version}: Semantic version of agent code
  \item \texttt{service\_name}: Hosting service identifier
  \item \texttt{deployment\_id}: Specific deployment instance
  \item \texttt{model\_id}: Underlying model identifier
  \item \texttt{trust\_level}: Autonomy level (L1--L4)
\end{itemize}

\subsection{Request identity fields}

\begin{itemize}
  \item \texttt{request\_id}: Unique request identifier
  \item \texttt{correlation\_id}: Correlation ID for distributed tracing
  \item \texttt{user\_id}: Requesting user (if applicable)
  \item \texttt{tenant\_id}: Tenant context
  \item \texttt{timestamp}: Request timestamp
\end{itemize}

\section{Audit reservation pattern}

The \textbf{reservation-before-action} pattern ensures audit completeness:

\begin{enumerate}
  \item \textbf{Reserve}: Before executing an action, reserve an audit event
        with status ``reserved''
  \item \textbf{Execute}: Perform the action
  \item \textbf{Complete}: Update the audit event with outcome (success/failure)
\end{enumerate}

If the action crashes without completing the audit, the reservation becomes
an ``orphan'' that indicates a governance failure requiring investigation.

\begin{tcolorbox}[colback=sgamber!10,colframe=sgamber!70,title=\faIcon{exclamation-triangle} Critical Requirement]
No irreversible action may execute without a successful audit reservation.
If the reservation fails, the action must be blocked.
\end{tcolorbox}

\section{Alert severity classes}

\begin{table}[htbp]
\centering
\caption{Alert Severity Classes}
\label{tab:alert-classes}
\begin{tabularx}{\textwidth}{@{}l l X@{}}
\toprule
\textbf{Severity} & \textbf{Response Time} & \textbf{Example Conditions} \\
\midrule
Critical & $< 15$ minutes & Identity missing, audit failure, circuit breaker open \\
High & $< 1$ hour & Orphaned audit, topic drift, latency spike \\
Medium & $< 4$ hours & Metric degradation, approval bottlenecks \\
Informational & Next business day & Successful operations, recovered services \\
\bottomrule
\end{tabularx}
\end{table}

\section{Condition codes}

Condition codes provide machine-readable alert classification:

\begin{table}[htbp]
\centering
\footnotesize
\caption{Governance Condition Codes}
\label{tab:condition-codes}
\begin{tabularx}{\textwidth}{@{}l l l X@{}}
\toprule
\textbf{Code} & \textbf{Severity} & \textbf{Control} & \textbf{Description} \\
\midrule
REG-ING-001 & Critical & AGT-02 & Tool called not in allow-list \\
REG-ING-002 & Critical & AGT-03 & Identity envelope missing \\
REG-AUD-001 & Critical & MON-01 & Audit reservation failed \\
REG-AUD-002 & High & MON-01 & Orphaned audit event detected \\
REG-SCHEMA-001 & High & MON-01 & Audit payload schema violation \\
REG-DRIFT-001 & High & EVAL-01 & Topic similarity below threshold \\
REG-LATENCY-001 & High & MON-02 & P95 latency exceeded threshold \\
REG-REFUSAL-001 & Medium & EVAL-01 & Refusal rate elevated \\
REG-CB-001 & Critical & MON-02 & Circuit breaker opened \\
\bottomrule
\end{tabularx}
\end{table}

\section{Circuit-break rules}

Circuit breakers protect against cascade failures:

\begin{table}[htbp]
\centering
\caption{Circuit-Break Triggers}
\label{tab:circuit-break}
\begin{tabularx}{\textwidth}{@{}X l l@{}}
\toprule
\textbf{Trigger Condition} & \textbf{Duration} & \textbf{Recovery} \\
\midrule
3 consecutive critical alerts in 15 minutes & Open until manual close & Human verification required \\
5 consecutive failed health checks & 30 seconds auto-retry & Auto-close on success \\
$>$ 10\% error rate in 5-minute window & Open until manual close & Investigation required \\
\bottomrule
\end{tabularx}
\end{table}

\section{Drift and anomaly detection}

\subsection{Population Stability Index (PSI)}

For detecting distribution drift in model inputs/outputs:

\begin{equation}
\boxed{
PSI = \sum_{i=1}^{k} (A_i - E_i) \cdot \ln\left(\frac{A_i}{E_i}\right)
}
\label{eq:psi}
\end{equation}

where $A_i$ is actual proportion in bin $i$ and $E_i$ is expected proportion.

\begin{table}[htbp]
\centering
\caption{PSI Interpretation}
\begin{tabular}{l l}
\toprule
\textbf{PSI Range} & \textbf{Interpretation} \\
\midrule
$PSI < 0.10$ & No significant drift \\
$0.10 \leq PSI < 0.25$ & Moderate drift, investigate \\
$PSI \geq 0.25$ & Significant drift, action required \\
\bottomrule
\end{tabular}
\end{table}

\subsection{Robust Z-Score}

For anomaly detection with outlier resistance:

\begin{equation}
\boxed{
Z_{\text{robust}} = \frac{x - \text{median}(X)}{1.4826 \times \text{MAD}(X)}
}
\label{eq:robust-z}
\end{equation}

where MAD is the Median Absolute Deviation. The factor 1.4826 normalizes MAD
to be consistent with standard deviation for normal distributions.

\begin{table}[htbp]
\centering
\caption{Z-Score Thresholds}
\begin{tabular}{l l}
\toprule
\textbf{Z-Score} & \textbf{Action} \\
\midrule
$|Z| < 2$ & Normal \\
$2 \leq |Z| < 3$ & Monitor \\
$|Z| \geq 3$ & Alert \\
\bottomrule
\end{tabular}
\end{table}

\section{Required alert payload}

Every alert must include:

\begin{table}[htbp]
\centering
\footnotesize
\caption{Required Alert Payload Fields}
\begin{tabularx}{\textwidth}{@{}l l X@{}}
\toprule
\textbf{Field} & \textbf{Type} & \textbf{Description} \\
\midrule
\texttt{request\_id} & string & Request that triggered the alert \\
\texttt{correlation\_id} & string & Correlation for distributed tracing \\
\texttt{agent\_id} & string & Agent that generated the condition \\
\texttt{severity} & enum & Critical, High, Medium, Informational \\
\texttt{condition\_code} & string & Machine-readable condition code \\
\texttt{condition\_message} & string & Human-readable description \\
\texttt{evidence\_source} & string & Path to supporting evidence \\
\texttt{observed\_value} & string & The value that triggered the alert \\
\texttt{threshold\_value} & string & The threshold that was violated \\
\texttt{recommended\_action} & string & Suggested remediation \\
\texttt{timestamp} & datetime & When the condition was detected \\
\bottomrule
\end{tabularx}
\end{table}

\section{Operator triage checklist}

When an alert is received:

\begin{enumerate}
  \item Confirm the request\_id and correlation\_id
  \item Confirm the agent\_id and deployment context
  \item Confirm the condition\_code and severity
  \item Verify the observed value against the threshold
  \item Check the evidence source for additional context
  \item Determine if this is a known issue (check Known Issues Registry)
  \item Execute recommended action or escalate
  \item Document resolution in incident ticket
\end{enumerate}

\section{Minimum evidence to close an alert}

\begin{itemize}
  \item Root cause identified
  \item Mitigation applied
  \item Audit trail shows successful processing post-fix
  \item No recurrence within observation window (15 minutes minimum)
\end{itemize}